\chapter[致谢]{致\quad 谢}\chaptermark{致\quad 谢}% syntax: \chapter[目录]{标题}\chaptermark{页眉}
\thispagestyle{noheaderstyle}% 如果需要移除当前页的页眉
%\pagestyle{noheaderstyle}% 如果需要移除整章的页眉

时间飞逝，很快就到了硕士毕业的时候了。借此论文完成的机会，在此向多年来给我关心和帮组的老师、同学、家人和朋友们表示衷心的感谢！

首先，要感谢我的父母，父母为我付出了毕生的精力，将我抚养成人。在我最困难的时候帮助我，虽然他们都不是科研人员，也没有很高的学术水平。但无论时科研上的烦心事，还是生活上的琐事，都会认真倾听，帮我解决问题。即使解决不了也会开导我。

接着向我敬爱的小导师——郑超老师表达我的感谢之情。郑超老师严谨认真、学识渊博、思维敏捷且富有激情。在我硕士在读期间，郑超老师对我的研究生课程学习、课题研究等方面给出了很多的指导和帮助，让我开阔了思路，培养了独立进行研究和工程设计的思维和基本素养。

其次，我要感谢组长和我的导师刘庆云老师在研究生阶段对我的帮助。刘老师视野开阔、知识渊博，并且十分关心学生的成长。每次都能发现学生工作、学习中的问题，并提出十分有指导意义的意见。

感谢第二研究室处理架构组的周舟老师、杨嵘老师、朱宇佳老师、李佳老师、李钊老师、杨威老师等。在实验室里老师们对每一个学生都提供了真心的帮助，在科研学习上，能得到他们的帮助，让我受益颇多。我也会学习各位老师对科研锐意进取、对生活充满热情的学习生活态度，在今后的工作生活中继续进步。

感谢在实验室给了我很大帮助的师兄师姐们：陆秋文、张曦、谭国林、黄彩云，他们在我最迷茫无助的时候给了我很多的帮助。特别是陆秋文师兄，在工作压力很大的情况下，对实验室的每个师弟师妹都给与了很多帮助，看到的有价值的知识、文章也第一时间分享给师弟师妹们，给我树立了优秀的榜样。感谢组里的崔一鸣、罗成志、李白杨师弟，邹雨婷、柳幼婷、周钊宇、王钰、尚晶晶、马小青师妹，在日常学习和生活中的陪伴。

还更要感谢同一届的李杰、贺文梁、张中一、杨兴东、张晓欧、田湘、钟友兵同学，在三年研究生生活中他们是陪伴我最多的人，虽然大家的研究方向各不相同，但一直在一起拼搏、一起奋斗，总找机会一起分享学习上的经验和成果。

\begin{flushright}
陈冠林

2019年6月写于国科大
\end{flushright}
\chapter{作者简历及攻读学位期间发表的学术论文与研究成果}

\section*{作者简历}

\subsection*{本文作者}
陈冠林，安徽省合肥市人，中共预备党员，中国科学院信息工程研究所硕士研究生。

2012.9——2016.6 北京理工大学~~软件学院

2016.9——2019.6 中国科学院信息工程研究所
\subsection*{完成项目}
去重降载专项

\section*{指导老师简介}
刘庆云，男，1980年生，博士，ACM和IEEE会员，中国计算机学会会员，2010年获得新疆科技进步一等奖，2012年入选信息工程研究所(下简称信工所)首届青年之星，2013年获得信工所重大科技进展奖。现任信工所高级工程师，硕士生导师，信息智能处理技术研究室（第二研究室）处理架构组组长，主要研究方向为信息内容安全，侧重网络安全系统体系架构、网络内容深度分析与检测理论与技术。作为主要负责人，先后承担国家863计划重大项目、国家242信息安全计划项目以及工信部、发改委等多项国家级重大信息安全科研项目。主持研发的多个系统在国家重要部门实际应用效果显著。近三年发表论文20篇，申请专利12篇，软著6项。
\cleardoublepage[plain]% 让文档总是结束于偶数页，可根据需要设定页眉页脚样式，如 [noheaderstyle]

